\chapter{संख्यात्मक श्रेणी}\label{ch:b1:series}

अनंत संख्याओं को जोड़ने का अर्थ है आंशिक योगों की सीमा लेना --- न इससे अधिक,
न कम। यह अध्याय परिभाषाएँ और प्रथम वर्ष में प्रयोग योग्य अभिसरण-कसौटियाँ
खड़ी करता है: धनात्मक पदों के लिए तुलना और तुल्यताएँ, \omterm{thm:b1:series:ratio}{अनुपात कसौटी}, \omterm{thm:b1:series:riemann}{रीमान
श्रेणी} देने वाली समाकल-तुलना, \omterm{thm:b1:series:absolute}{निरपेक्ष अभिसरण}, और \omterm{thm:b1:series:alternating}{एकांतरित श्रेणी} प्रमेय।
सूक्ष्मतर सिद्धांत (श्रेणियों के गुणनफल, पैकेटों में योग, फलनों की \omterm{def:b1:series:def}{श्रेणी})
द्वितीय वर्ष का है।

\section{सामान्य बातें}

\begin{definition}\label{def:b1:series:def}
अनुक्रम $(u_n)$ दिया हो, तो \emph{श्रेणी}\index{श्रेणी} $\sum u_n$
\emph{आंशिक योगों} $S_N = \sum_{n=0}^{N} u_n$ का अनुक्रम है। $(S_N)$ के
अभिसरित होने पर श्रेणी \emph{अभिसरित} होती है; सीमा \emph{योग}
$\sum_{n=0}^{\infty} u_n$ है, और $R_N = \sum_{n > N} u_n = S - S_N$
\emph{शेषफल} है, जो $0$ की ओर जाता है।
\end{definition}

\begin{example}[गुणोत्तर श्रेणी]\label{ex:b1:series:geometric}
$q \in \C$ के लिए: $\;S_N = \sum_{n=0}^{N} q^n = \frac{1 - q^{N+1}}{1-q}$
($q \neq 1$)। \omterm{def:b1:series:def}{श्रेणी} अभिसरित होती है तभी जब $\abs q < 1$
(\cref{exo:b1:seq:3}), जहाँ\index{गुणोत्तर श्रेणी}
\[
\sum_{n=0}^{\infty} q^n = \frac{1}{1 - q} .
\]
\end{example}

\begin{example}[आवर्ती दशमलव गुणोत्तर श्रेणी हैं]\label{ex:b1:series:periodicdecimal}
$0.363636\dots$ कौन-सी संख्या है? उसका लेखन ही एक \omterm{def:b1:series:def}{श्रेणी} है:
\[
0.\overline{36} = \sum_{k=1}^{\infty} \frac{36}{100^k}
= 36\cdot\frac{1/100}{1 - 1/100} = \frac{36}{99} =
\frac{4}{11} ,
\]
जो $q = \frac{1}{100}$ वाले गुणोत्तर योग से मिलता है। व्यापक रूप से, $p$
अंकों का सदा दोहराता खंड $B$ $\frac{B}{10^p - 1}$ के बराबर है --- यही
क्रियाविधि \cref{pb:b1:reals:1} की आवर्तिता कसौटी के पीछे थी, जिसे इस अध्याय
की भाषा अंततः एक पंक्ति में कह देती है: \omterm{pb:b1:reals:1}{दशमलव प्रसार} एक अभिसारी \omterm{def:b1:series:def}{श्रेणी} है,
जो ठीक तब अंततः आवर्ती होती है जब उसका योग परिमेय हो। अध्याय 10 की
अंक-मशीनरी, जो वहाँ नंगे उच्चतमों से बनी थी, वेश बदलकर चलती श्रेणी-थ्योरी ही
थी।
\end{example}

\begin{proposition}[पहले तथ्य]\label{prop:b1:series:first}
\begin{enumerate}
  \item यदि $\sum u_n$ अभिसरित होती है, तो $u_n \to 0$। (विलोम \emph{असत्य}
  है: हरात्मक \omterm{def:b1:series:def}{श्रेणी}।)
  \item रैखिकता: अभिसारी श्रेणियाँ अपेक्षित योगों के साथ जुड़ती और मापित
  होती हैं।
  \item (दूरबीनी\index{दूरबीनी श्रेणी}) $\sum (v_{n+1} - v_n)$ अभिसरित होती
  है तभी जब $(v_n)$ अभिसरित हो, और योग $\lim v_n - v_0$।
  \item परिमित कितने पद बदलने से अभिसरण नहीं बदलता (केवल योग बदलता है)।
\end{enumerate}
\end{proposition}

\begin{proof}
(1) $u_N = S_N - S_{N-1} \to S - S = 0$। हरात्मक \omterm{def:b1:series:def}{श्रेणी} में $u_n = \frac1n
\to 0$ है, फिर भी वह अपसरित होती है (\cref{exo:b1:seq:5})। (2) सीमाओं पर
संक्रियाएँ। (3) $S_N = v_{N+1} - v_0$। (4) आंशिक योग अंततः अचर राशि भर बदल
जाते हैं।
\end{proof}

\begin{example}[गुणोत्तर शेषफल से अंकों की योजना]\label{ex:b1:series:georemainder}
$\abs q < 1$ के लिए \omterm{ex:b1:series:geometric}{गुणोत्तर श्रेणी} का शेषफल स्पष्ट है:
\[
R_N = \sum_{n = N+1}^{\infty} q^n = \frac{q^{N+1}}{1 - q} .
\]
इससे यथार्थता के लक्ष्य किसी भी संगणना से पहले पद-संख्या में बदल जाते हैं।
$\sum_{n\geq0} \bigl(\frac13\bigr)^n = \frac32$ का मान $10^{-10}$ के भीतर
निकालने के लिए: $\frac{(1/3)^{N+1}}{2/3} \leq 10^{-10}$ चाहिए, अर्थात्
$3^{N} \geq \frac{3}{2}\cdot 10^{10}$, अर्थात् $N \geq 22$ (क्योंकि $3^{22}
\approx 3.1\cdot10^{10}$): अतः तेईस पद, और वह भी पहले से ज्ञात। सप्ताहांत
समस्याओं का हर गुणोत्तर-दर वाला आकलन ($\ln 2$ के लिए $\frac13$-श्रेणी,
\cref{pb:b1:taylor:1} में मैकिन की प्रतिलोम स्पर्शज्याएँ) यही दो पंक्ति का
बजट है, बस पेशेवर पोशाक में।
\end{example}

\begin{example}[एक लंबी दूरबीन]\label{ex:b1:series:telescope3}
$\sum_{n\geq1} \frac{1}{n(n+1)(n+2)}$ संगणित कीजिए। आंशिक भिन्नें
(\cref{ch:b1:fractions}):
\[
\frac{1}{n(n+1)(n+2)}
= \frac{1/2}{n} - \frac{1}{n+1} + \frac{1/2}{n+2}
= \frac12\Bigl(\frac{1}{n(n+1)} - \frac{1}{(n+1)(n+2)}\Bigr),
\]
जहाँ दूसरा रूप --- $w_n = \frac{1}{n(n+1)}$ के क्रमागत मानों का अंतर ---
दूरबीनी रूप है। अतः
\[
\sum_{n=1}^{N} \frac{1}{n(n+1)(n+2)}
= \frac12\Bigl(w_1 - w_{N+1}\Bigr)
= \frac12\Bigl(\frac12 - \frac{1}{(N+1)(N+2)}\Bigr)
\longrightarrow \frac14 .
\]
समापन का सार: तीन-पद वाली आंशिक भिन्नें जैसी लिखी हों वैसी शायद ही दूरबीनी
होती हैं; पहले उन्हें किसी अंतर $w_n - w_{n+1}$ में पुनर्समूहित कीजिए --- और
पुरस्कार केवल अभिसरण नहीं, बल्कि यथार्थ योग है, जो कोई तुलना-कसौटी कभी नहीं
देती।
\end{example}

\section{ऋणेतर पदों वाली श्रेणी}

\begin{theorem}[परिबद्ध आंशिक योग]\label{thm:b1:series:positive}
यदि सभी $n$ के लिए $u_n \geq 0$, तो आंशिक योग बढ़ते हैं, अतः: $\sum u_n$
अभिसरित होती है $\iff$ उसके आंशिक योग ऊपर से परिबद्ध हैं। इससे \emph{तुलना
कसौटी} मिलती है: यदि सभी (बड़े) $n$ के लिए $0 \leq u_n \leq v_n$, तो
\begin{align*}
\sum v_n \text{ अभिसरित होती है} &\implies \sum u_n \text{ अभिसरित होती है},\\
\sum u_n \text{ अपसरित होती है} &\implies \sum v_n \text{ अपसरित होती है}.
\end{align*}
और \emph{तुल्यता कसौटी}: यदि $v_n \geq 0$ के साथ $u_n \sim v_n$, तो दोनों
श्रेणियों की प्रकृति एक ही है।
\end{theorem}

\begin{proof}
पहले बिंदु के लिए एकदिष्ट सीमा प्रमेय (\cref{thm:b1:seq:monotone}); और दूसरे
के लिए आंशिक योगों की तुलना। तुल्यताएँ: बड़े $n$ के लिए $\frac12 v_n \leq
u_n \leq 2 v_n$ ($\varepsilon = \frac12$ के साथ $\sim$ की परिभाषा), और तुलना
दोनों ओर लागू हो जाती है।
\end{proof}

\begin{example}[एक तुल्यता जो अपसरण सिद्ध कर देती है]\label{ex:b1:series:equivdiverge}
$\sum_{n\geq1} n\sin\dfrac{1}{n^2}$ की प्रकृति? चूँकि $0$ पर $\frac{1}{n^2}
\to 0$ और $\sin h \sim h$:
\[
n\sin\frac{1}{n^2} \;\sim\; n\cdot\frac{1}{n^2} = \frac1n ,
\]
और तुल्यता कसौटी हरात्मक \omterm{def:b1:series:def}{श्रेणी} का अपसरण स्थानांतरित कर देती है: अतः अपसारी
--- यद्यपि पद $0$ की ओर जाते हैं। एक प्रसार, एक पैमाना, एक निर्णय; और यही
दो-पदीय प्रतिरूप (तुल्यता, फिर रीमान या गुणोत्तर की खोज)
\cref{exo:b1:series:3} की चारों श्रेणियों का निर्णय कर देता है।
\end{example}

\begin{example}[एक ही पंक्ति में तुल्यता कसौटी]\label{ex:b1:series:equivtest}
$\sum_{n \geq 1} \frac{\sqrt{n+1} - \sqrt n}{n}$ की प्रकृति? अंश का संयुग्मी
लीजिए:
\[
\frac{\sqrt{n+1} - \sqrt n}{n}
= \frac{1}{n\,(\sqrt{n+1} + \sqrt n)}
\sim \frac{1}{2\,n^{3/2}} ,
\]
यह अभिसारी रीमान पैमाना ($\alpha = \frac32 > 1$) है: अतः \omterm{def:b1:series:def}{श्रेणी} अभिसरित होती
है। पूरा निर्णय एक तुल्यता और एक खोज में हो गया --- बशर्ते पद ऋणेतर हों, जो
वे हैं। समापन का सार: धनात्मक श्रेणियों के लिए पूरा अभिसरण-सिद्धांत
\emph{पैमानों का शब्दकोश} ($n^{-\alpha}$, $q^n$, $\frac{1}{n(\ln
n)^\alpha}$) और पद को उसकी तुल्यता से बदलने की अनुमति भर है; विश्लेषणात्मक
काम अनंतस्पर्शी में है (\cref{ch:b1:taylor}), जोड़ने में कभी नहीं।
\end{example}

\begin{theorem}[समाकल-तुलना; रीमान श्रेणी]\label{thm:b1:series:riemann}
मान लीजिए $f$ $\intco{1}{+\infty}$ पर \omterm{def:b1:continuity:continuous}{संतत}, ऋणेतर और \emph{ह्रासमान} है। तब
\[
\int_1^{N+1} f(t)\,\dd t \;\leq\; \sum_{n=1}^{N} f(n) \;\leq\; f(1)
+ \int_1^{N} f(t)\,\dd t ,
\]
अतः $\sum f(n)$ अभिसरित होती है तभी जब $\bigl(\int_1^x f\bigr)$ परिबद्ध हो।
विशेष रूप से $\alpha \in \R$ के लिए:\index{रीमान श्रेणी}
\[
\sum_{n \geq 1} \frac{1}{n^\alpha} \text{ अभिसरित होती है}
\iff \alpha > 1,
\]
और $\sum_{n=1}^{N} \frac1n = \ln N + O(1)$।
\end{theorem}

\begin{proof}
$n \leq t \leq n+1$ के लिए एकदिष्टता $f(n+1) \leq f(t) \leq f(n)$ देती है;
$\intcc{n}{n+1}$ (जो $1$ लंबाई का खंड है) पर समाकलन करने पर:
\[
f(n+1) \;\leq\; \int_n^{n+1} f(t)\,\dd t \;\leq\; f(n) .
\]
$n = 1, \dots, N-1$ के लिए दाईं असमिकाओं का योग $\int_1^{N} f \leq
\sum_{n=1}^{N-1} f(n)$ देता है, जिससे $f(N) \leq f(1)$ जोड़ने पर ऊपरी घेराव
मिलता है; और $n = 1, \dots, N$ के लिए बाईं असमिकाओं का योग $\sum_{n=2}^{N+1}
f(n) \leq \int_1^{N+1} f$ देता है, जो पुनःसूचीबद्ध करने पर निचला घेराव है।
अभिसरण: आंशिक योग और \omterm{thm:b1:integration:def}{समाकल} $\int_1^x f$ एक-दूसरे को अचर $f(1)$ के भीतर
परिबद्ध करते हैं, और दोनों अह्रासमान हैं, अतः एक परिबद्ध है तभी जब दूसरा हो
(\cref{thm:b1:series:positive})। $f(t) = t^{-\alpha}$ ($\alpha \neq 1$) के
लिए: $\int_1^x t^{-\alpha}\dd t = \frac{x^{1-\alpha} - 1}{1 - \alpha}$, जो
$\alpha > 1$ होने पर परिबद्ध है; $\alpha = 1$ के लिए \omterm{thm:b1:integration:def}{समाकल} $\ln x \to
\infty$ है, और घेराव $\ln(N+1) \leq H_N \leq 1 + \ln N$ देता है। $\alpha
\leq 0$ के लिए पद $0$ की ओर जाते ही नहीं।
\end{proof}

\begin{example}[हरात्मक ढेर]\label{ex:b1:series:harmonicstack}
हरात्मक \omterm{def:b1:series:def}{श्रेणी} को $20$ पार करने के लिए कितने पद जोड़ने पड़ते हैं? घेराव
$\ln(N+1) \leq H_N \leq 1 + \ln N$ बिना कुछ जोड़े उत्तर दे देता है: $H_N
\geq 20$ के लिए $1 + \ln N \geq 20$ चाहिए, अर्थात् $N \geq \eu^{19} \approx
1.8\cdot10^{8}$, और $\ln(N + 1) \geq 20$ होते ही, अर्थात् $N \approx
\eu^{20} \approx 4.9\cdot10^{8}$ होते ही, वह गारंटीशुदा है। (सप्ताहांत
समस्या इसे ऑयलर के अचर के द्वारा $N \approx \eu^{20 - \gamma} \approx
2.7\cdot10^{8}$ तक तीक्ष्ण कर देती है।) समापन का सार: समाकल-तुलना केवल
अभिसरण तय नहीं करती --- वह आंशिक योगों को लघुगणकीय यथार्थता के साथ
\emph{ढूँढ़} भी देती है, और एक निराशाजनक संगणना (करोड़ों पद) को दो पंक्ति के
आकलन में बदल देती है।
\end{example}

\begin{theorem}[अनुपात कसौटी (दालंबेर)]\label{thm:b1:series:ratio}
मान लीजिए $\frac{u_{n+1}}{u_n} \to \ell$ के साथ $u_n > 0$।\index{अनुपात
कसौटी}
\begin{itemize}
  \item यदि $\ell < 1$: $\sum u_n$ अभिसरित होती है;
  \item यदि $\ell > 1$: $u_n \to +\infty$, अतः अपसरण;
  \item यदि $\ell = 1$: कोई निष्कर्ष नहीं ($\sum \frac1n$ अपसरित होती है,
  $\sum \frac{1}{n^2}$ अभिसरित)।
\end{itemize}
\end{theorem}

\begin{proof}
यदि $\ell < 1$, तो $q \in \intoo{\ell}{1}$ नियत कीजिए: किसी $N$ के आगे
$u_{n+1} \leq q\,u_n$, अतः आगमन से $u_n \leq u_N q^{\,n-N}$: अर्थात् किसी
\omterm{ex:b1:series:geometric}{गुणोत्तर श्रेणी} से तुलना। यदि $\ell > 1$: तो किसी $N$ के आगे अनुक्रम $(u_n)$
वर्धमान है, अतः वह $0$ की ओर नहीं जा सकता (उसकी सीमा, यदि हो, $\geq u_N > 0$
है); और \cref{prop:b1:series:first}~(1) से अपसरण --- वस्तुतः $q > 1$ के साथ
$u_n \geq u_N q^{n-N}$ $u_n \to \infty$ दे देता है।
\end{proof}

\begin{example}\label{ex:b1:series:ratio}
$\sum \frac{x^n}{n!}$ प्रत्येक $x > 0$ के लिए अभिसरित होती है: अनुपात
$\frac{x}{n+1} \to 0$। उसका योग $\eu^x$ है: $\intcc{0}{x}$ पर टेलर--लाग्रांज
(\cref{thm:b1:taylor:lagrange}) से,
\[
\Bigl| \eu^x - \sum_{k=0}^{n} \frac{x^k}{k!} \Bigr|
\leq \eu^{x}\, \frac{x^{n+1}}{(n+1)!} \xrightarrow[n\to\infty]{} 0 ,
\]
जहाँ परिबंध $0$ की ओर जाता है, क्योंकि क्रमगुणित हावी है
(\cref{exo:b1:integration:9}~(1) ने भी वही तथ्य प्रयुक्त किया था)। वही तर्क
पूरे $\R$ पर $\sin$, $\cos$, $\sinh$, $\cosh$ की श्रेणियों का योग निकाल देता
है।
\end{example}

\begin{example}[अनुपात कसौटी पर्याप्त है, आवश्यक नहीं]\label{ex:b1:series:rationotnecessary}
सम $n$ के लिए $u_n = 2^{-n}$ और विषम $n$ के लिए $u_n = 2^{-n-2}$ लीजिए।
क्रमागत अनुपात $\frac{1}{8}$ और $\frac12\cdot4 = 2$ के बीच दोलन करते हैं,
अतः $\frac{u_{n+1}}{u_n}$ की कोई सीमा नहीं है और दालंबेर मौन है --- फिर भी
$u_n \leq 2^{-n}$ है और तुलना कसौटी अभिसरण तुरंत तय कर देती है। इस कसौटी की
परिकल्पना (अनुपात का \emph{अभिसरित होना}) सचमुच एक प्रतिबंध है: वह उन पदों
पर सूट करती है जिनकी एक ही प्रभावी गुणात्मक संरचना हो (क्रमगुणित, घातें), और
जो कुछ भी साँस लेता है उस पर विफल हो जाती है। जब अनुपात बदमाशी करें, तो पीछे
हटकर किसी गुणोत्तर अन्वालोप से तुलना कीजिए --- और \omterm{thm:b1:series:ratio}{अनुपात कसौटी} कभी इससे अधिक
थी भी नहीं, जैसा उसकी उपपत्ति दिखाती है।
\end{example}

\begin{example}[क्रमगुणितों की लड़ाई पर अनुपात कसौटी]\label{ex:b1:series:ratiobinom}
$\sum_{n\geq0} \dfrac{(n!)^2}{(2n)!}$ की प्रकृति (केंद्रीय द्विपद गुणांकों
के व्युत्क्रम, गुणक $n + 1$ तक)? अनुपात क्रमगुणितों को ढहा देता है:
\[
\frac{u_{n+1}}{u_n}
= \frac{((n+1)!)^2}{(n!)^2}\cdot\frac{(2n)!}{(2n+2)!}
= \frac{(n+1)^2}{(2n+1)(2n+2)}
\longrightarrow \frac14 < 1 :
\]
अतः अभिसारी, और वह भी काफ़ी अंतर से --- पद मूलतः $4^{-n}$ की तरह क्षय होते
हैं, जो \cref{pb:b1:integration:1} के $\binom{2n}{n} \geq \frac{4^n}{2n+1}$
के अनुरूप है। समापन का सार: क्रमगुणितों के भागफल ठीक वही हैं जिन्हें \omterm{thm:b1:series:ratio}{अनुपात
कसौटी} पचा लेती है --- हर क्रमगुणित $n$ के किसी परिमेय फलन में कट जाता है,
जिसकी सीमा अग्र पदों से पढ़ ली जाती है।
\end{example}

\section{निरपेक्ष अभिसरण; एकांतरित श्रेणी}

\begin{theorem}[निरपेक्ष अभिसरण]\label{thm:b1:series:absolute}
यदि $\sum \abs{u_n}$ अभिसरित होती है (\emph{निरपेक्ष अभिसरण}\index{निरपेक्ष
अभिसरण}), तो $\sum u_n$ भी अभिसरित होती है, और $\bigl|\sum u_n\bigr| \leq
\sum \abs{u_n}$। यह वास्तविक और सम्मिश्र दोनों प्रकार के पदों के लिए सत्य
है।
\end{theorem}

\begin{proof}
आंशिक योग $M > N$ के लिए यह संतुष्ट करते हैं (कोशी कसौटी,
\cref{thm:b1:seq:complete}):
\[
\abs{S_M - S_N} = \Bigl| \sum_{n=N+1}^{M} u_n \Bigr|
\leq \sum_{n=N+1}^{M} \abs{u_n},
\]
जो बड़े $N$ के लिए छोटा है, क्योंकि $\sum\abs{u_n}$ के आंशिक योग \omterm{def:b1:seq:cauchy}{कोशी
अनुक्रम} बनाते हैं। अतः $(S_N)$ कोशी है, इसलिए अभिसारी। और असमिका परिमित
त्रिभुज असमिका से सीमा तक चली जाती है।
\end{proof}

\begin{example}[निरपेक्ष अभिसरण, वास्तविक और सम्मिश्र]\label{ex:b1:series:absexamples}
$\sum_{n\geq1} \frac{\sin n}{n^2}$: पदों के चिह्न अनियमित रूप से बदलते हैं
(वस्तुतः $(\sin n)$ $\intcc{-1}{1}$ में \omterm{def:b1:topology:dense}{सघन} है, \cref{exo:b1:seq:12}), और
कोई एकांतरित संरचना दिखाई नहीं देती। \omterm{thm:b1:series:absolute}{निरपेक्ष अभिसरण} एक ही झटके में सब कुछ
बचा लेता है: $\bigl|\frac{\sin n}{n^2}\bigr| \leq \frac{1}{n^2}$, जो अभिसारी
पैमाना है, अतः \omterm{def:b1:series:def}{श्रेणी} अभिसरित होती है। वही ढाल $\C$ पर भी काम करती है:
$\sum_{n\geq1}\frac{\eu^{\iu n}}{n^2}$ अभिसरित होती है क्योंकि
$\bigl|\frac{\eu^{\iu n}}{n^2}\bigr| = \frac{1}{n^2}$ --- अर्थात् \omterm{def:b1:complex:field}{मापांक}
योग्य होते ही चिह्नों के प्रतिरूप, चाहे वे द्विविमीय ही क्यों न हों, निरर्थक
हो जाते हैं। समापन का सार: इस अध्याय का एकमात्र औज़ार \omterm{thm:b1:series:absolute}{निरपेक्ष अभिसरण} है जो
कभी यह नहीं पूछता कि चिह्न कैसे व्यवस्थित हैं; उसे पहले आज़माइए
(\cref{met:b1:series:decide}), और नाज़ुक कसौटियाँ उन्हीं श्रेणियों के लिए
रखिए जो उसमें विफल हो जाएँ।
\end{example}

\begin{theorem}[एकांतरित श्रेणी कसौटी]\label{thm:b1:series:alternating}
मान लीजिए $(a_n)$ ह्रासमान है और $a_n \to 0$। तब \omterm{thm:b1:series:alternating}{एकांतरित श्रेणी} $\sum
(-1)^n a_n$ अभिसरित होती है; उसका योग किन्हीं दो क्रमागत आंशिक योगों के बीच
रहता है, और\index{एकांतरित श्रेणी}
\[
\abs{R_N} = \Bigl| \sum_{n > N} (-1)^n a_n \Bigr| \leq a_{N+1} .
\]
\end{theorem}

\begin{proof}
सम और विषम आंशिक योग संलग्न हैं: $S_{2p+2} - S_{2p} = a_{2p+2} - a_{2p+1}
\leq 0$ (ह्रासमान), $S_{2p+1} - S_{2p-1} = a_{2p} - a_{2p+1} \geq 0$
(वर्धमान), और $S_{2p} - S_{2p+1} = a_{2p+1} \to 0$।
\cref{thm:b1:seq:adjacent} से उनकी उभयनिष्ठ सीमा $S$ है, जिसे दो उपानुक्रमों
की कसौटी (\cref{prop:b1:seq:subsequences}) $(S_N)$ की सीमा बना देती है; इसके
अतिरिक्त $S$ क्रमागत आंशिक योगों के बीच फँसा है, और $\abs{S - S_N}$ अधिक से
अधिक अगले तक के अंतर $a_{N+1}$ जितना है।
\end{proof}

\begin{example}[एकांतरित हरात्मक श्रेणी]\label{ex:b1:series:ln2}
$\sum_{n \geq 1} \frac{(-1)^{n-1}}{n}$ अभिसरित होती है (एकांतरित कसौटी) पर
निरपेक्ष रूप से नहीं (हरात्मक \omterm{def:b1:series:def}{श्रेणी})। उसका योग $\ln 2$ है: परिमित गुणोत्तर
सर्वसमिका $\frac{1}{1+t} = \sum_{k=0}^{n-1} (-t)^k + \frac{(-t)^n}{1+t}$ का
$\intcc{0}{1}$ पर समाकलन कीजिए:
\[
\ln 2 = \sum_{k=1}^{n} \frac{(-1)^{k-1}}{k}
+ (-1)^n \int_0^1 \frac{t^n}{1+t}\,\dd t ,
\qquad
0 \leq \int_0^1 \frac{t^n}{1+t}\,\dd t \leq \frac{1}{n+1} \to 0 .
\]
अभिसरण पीड़ादायक रूप से धीमा है ($R_N \approx \frac{1}{N}$) --- \omterm{thm:b1:series:alternating}{एकांतरित
श्रेणी} निरसन से अभिसरित होती हैं, लघुता से नहीं।
\end{example}

\begin{omfigure}
\begin{tikzpicture}[x=0.78cm, y=7.2cm]
  \draw[-Stealth] (0.3,0.42) -- (11,0.42)
    node[below, font=\small] {$N$};
  \draw (1,0.415) -- (1,0.425) node[below=1pt, font=\small]
    {$1$};
  \draw (5,0.415) -- (5,0.425) node[below=1pt, font=\small]
    {$5$};
  \draw (10,0.415) -- (10,0.425) node[below=1pt, font=\small]
    {$10$};
  \draw[gray, dashed] (0.3,0.6931) -- (10.7,0.6931)
    node[right, font=\small, text=black] {$\ln 2$};
  \node[left, font=\small] at (0.3,1) {$1$};
  \node[left, font=\small] at (0.3,0.5) {$0.5$};
  \draw[omDef, thick]
    plot coordinates {(1,1) (2,0.5) (3,0.8333) (4,0.5833)
    (5,0.7833) (6,0.6167) (7,0.7595) (8,0.6345)
    (9,0.7456) (10,0.6456)};
  \fill[omDef] (1,1) circle (1.5pt);
  \fill[omDef] (2,0.5) circle (1.5pt);
  \fill[omDef] (3,0.8333) circle (1.5pt);
  \fill[omDef] (4,0.5833) circle (1.5pt);
  \fill[omDef] (5,0.7833) circle (1.5pt);
  \fill[omDef] (6,0.6167) circle (1.5pt);
  \fill[omDef] (7,0.7595) circle (1.5pt);
  \fill[omDef] (8,0.6345) circle (1.5pt);
  \fill[omDef] (9,0.7456) circle (1.5pt);
  \fill[omDef] (10,0.6456) circle (1.5pt);
\end{tikzpicture}

{\small एकांतरित हरात्मक \omterm{def:b1:series:def}{श्रेणी} $1 - \frac12 + \frac13 - \cdots$ के आंशिक
योग $S_N$ हर पद पर अपनी सीमा $\ln 2$ के ऊपर से कूदते हैं: विषम योग ऊपर से,
सम योग नीचे से, और हर कूद का आकार $\frac{1}{N+1}$। यह घेराव
\cref{thm:b1:series:alternating} की उपपत्ति को दृश्य बना देता है --- और
चिमटे का धीरे-धीरे बंद होना ($\abs{S_N - \ln 2} \approx \frac{1}{2N}$,
सप्ताहांत समस्या \cref{pb:b1:series:1}) ही वह कारण है कि $\ln 2$ कोई इस तरह
संगणित नहीं करता।}
\end{omfigure}

\begin{remark}[श्रेणियों की सामान्य भूलें]
(क) \emph{तुल्यता कसौटी को चिह्न चाहिए}: मान लीजिए $v_n =
\frac{(-1)^n}{\sqrt n}$ और $u_n = v_n + \frac1n$। तब $\frac{u_n}{v_n} = 1 +
\frac{(-1)^n}{\sqrt n} \to 1$, अतः $u_n \sim v_n$; फिर भी $\sum v_n$ अभिसरित
होती है (एकांतरित कसौटी) जबकि $\sum u_n = \sum v_n + \sum \frac1n$ अपसरित।
तुल्यता पदों के \emph{आकार} को नियंत्रित करती है, और चिह्नित श्रेणियों के
लिए आकार ही भाग्य नहीं है --- यह कसौटी (अंततः) ऋणेतर पदों के लिए ही कही गई
है और वहीं सत्य है। (ख) \emph{$u_n \to 0$ कुछ भी सिद्ध नहीं करता}: हरात्मक
\omterm{def:b1:series:def}{श्रेणी} शाश्वत प्रति-उदाहरण है; और विलोम दिशा
(\cref{prop:b1:series:first}~(1)) केवल एक त्वरित अपसरण-परीक्षण है। (ग)
\emph{अनुपात की सीमा $1$ मौन है, अभिसरण नहीं}: $\sum\frac1n$ और
$\sum\frac{1}{n^2}$ दोनों का अनुपात $\to 1$ है; रीमान पैमानों या समाकल-तुलना
पर जाइए। (घ) \emph{एकांतरित को ह्रासमान चाहिए}: $\sum \frac{(-1)^n}{n +
(-1)^n}$ एकांतरित लगती है पर उसे केवल प्रसार (\cref{exo:b1:series:5})
सँभालता है; और \cref{ch:b1:taylor} की सप्ताहांत समस्या (वहाँ प्रश्न 23)
दिखाती है कि एकदिष्टता के बिना कसौटी पूरी तरह विफल हो सकती है। (ङ)
\emph{समूहन और पुनर्क्रमण मुफ़्त नहीं हैं}: कोष्ठक डालना अभिसारी श्रेणियों
के लिए निरापद है पर अपसरण से अभिसरण गढ़ सकता है (जोड़ों में समूहबद्ध $1 - 1
+ 1 - \cdots$), और पुनर्क्रमण तो योग ही बदल सकता है --- यही नाटक इस अध्याय
की सप्ताहांत समस्या (\cref{pb:b1:series:1}) में मंचित है।
\end{remark}

\begin{method}[किसी श्रेणी की प्रकृति तय करना]\label{met:b1:series:decide}
\begin{enumerate}
  \item क्या $u_n \to 0$? यदि नहीं, तो अपसरण, और बात ख़त्म।
  \item ऋणेतर पद: $u_n$ की कोई तुल्यता खोजिए (प्रसार, \cref{ch:b1:taylor}!),
  और रीमान या गुणोत्तर पैमानों से तुलना कीजिए; क्रमगुणित और घातें \omterm{thm:b1:series:ratio}{अनुपात
  कसौटी} माँगती हैं; ह्रासमान $f(n)$ समाकल-तुलना माँगता है।
  \item चिह्न बदलते हों: पहले \omterm{thm:b1:series:absolute}{निरपेक्ष अभिसरण} आज़माइए; वह विफल हो तो
  एकांतरित कसौटी (\emph{ह्रासमान} होना सावधानी से जाँचिए); उससे आगे द्वितीय
  वर्ष के औज़ार।
\end{enumerate}
\end{method}

\begin{example}[विषम हर, आधी दूरबीन]\label{ex:b1:series:oddtelescope}
$\sum_{n \geq 1} \dfrac{1}{4n^2 - 1}$ संगणित कीजिए। आंशिक भिन्नें:
$\frac{1}{(2n-1)(2n+1)} = \frac12\bigl(\frac{1}{2n-1} -
\frac{1}{2n+1}\bigr)$, अतः
\[
\sum_{n=1}^{N} \frac{1}{4n^2 - 1}
= \frac12\Bigl(1 - \frac{1}{2N+1}\Bigr)
\longrightarrow \frac12 .
\]
$\sum \frac{1}{n(n+1)} = 1$ (\cref{exo:b1:series:1}) से तुलना कीजिए: वही
दूरबीनी ढाँचा, पर यहाँ क्रमागत पद विषम संख्याओं में दो-दो की दूरी पर हैं, और
गुणक $\frac12$ उस पद को दर्ज करता है। समापन का सार: दूरबीनी कोई युक्ति नहीं,
दृष्टिकोण का परिवर्तन है --- जब भी व्यापक पद किसी सीमा-युक्त अनुक्रम का अंतर
$w_n - w_{n+1}$ हो, तब योग $w_1 - \lim w$ होता है, ठीक
\cref{prop:b1:series:first}~(3) के अनुसार।
\end{example}

\begin{remark}[विश्लेषण की पाइपलाइन, पीछे मुड़कर]
इस अध्याय में इस खंड का विश्लेषण आ मिलता है, और हर कसौटी अपने पूर्वज का नाम
लेती है। परिबद्ध आंशिक योग एकदिष्ट सीमा प्रमेय (\cref{ch:b1:seq}) है, जो
स्वयं \cref{ch:b1:reals} का पूर्णता अभिगृहीत है; \omterm{thm:b1:series:absolute}{निरपेक्ष अभिसरण} कोशी कसौटी
है; \omterm{thm:b1:integration:def}{समाकल} कसौटी \cref{ch:b1:integration} का क्षेत्रफलों वाला घेराव है;
व्यापक पदों की तुल्यताएँ \cref{ch:b1:taylor} के प्रसार हैं; और एकांतरित
प्रमेय संलग्न-अनुक्रम प्रमेयिका अपने रविवारी वस्त्रों में। उलटी दिशा में
पढ़ने पर यह पाइपलाइन समझा देती है कि हर अध्याय \emph{किसलिए} था --- और उसमें
पिरोई सप्ताहांत समस्याएँ ($b$-आदिक अंक, चेज़ारो--स्टोल्ज़, अपरिमेयता-मशीनें,
\omterm{pb:b1:series:1}{ऑयलर का अचर}) वही कुछ विचार हैं जो निरंतर ऊँची होती ऊँचाई पर मिलते जाते हैं।
आगे आने वाला रैखिक बीजगणित विषय बदलता है, मानक नहीं: प्रमाणित त्रुटि के साथ
यथार्थ \omterm{def:b1:logic:statement}{कथन} की आदत सीमाओं से विमाओं तक की यात्रा में बची रहती है।
\end{remark}

\begin{remark}[श्रेणियाँ आगे कहाँ जाती हैं]
यह अध्याय इस खंड का विश्लेषण बंद करता है और तीन द्वार खोलता है। स्नातक वर्ष
2 के खंड में श्रेणियाँ एक चर पा लेती हैं ($\sum a_n x^n$: घात \omterm{def:b1:series:def}{श्रेणी}, अपनी
अभिसरण-त्रिज्या के साथ) और फिर फलन-मान वाला सिद्धांत (फूरिये \omterm{def:b1:series:def}{श्रेणी}); और
निरपेक्ष-बनाम-सप्रतिबंध अभिसरण का द्विभाजन, जिसे नीचे की सप्ताहांत समस्या
नाटकीय बना देती है, दोनों की आधारशिला बन जाता है। प्रायिकता में (स्नातक वर्ष
3 का खंड) विविक्त यादृच्छिक चरों की प्रत्याशाएँ श्रेणियाँ \emph{हैं}, और
\omterm{thm:b1:series:absolute}{निरपेक्ष अभिसरण} ही उन्हें सुपरिभाषित बनाता है। और \omterm{thm:b1:series:riemann}{रीमान श्रेणी} $\sum
n^{-s}$, सम्मिश्र $s$ तक बढ़ाकर, ज़ीटा फलन बन जाती है --- गणित की सबसे अधिक
अध्ययन की गई \omterm{def:b1:series:def}{श्रेणी}।
\end{remark}

\section{अभ्यास}

\begin{exercise}[$\star$]\label{exo:b1:series:1}
प्रकृति (और दूरबीनी होने पर योग भी) बताइए:
\[
\sum_{n\geq1} \frac{1}{n(n+1)},
\qquad
\sum_{n\geq2} \ln\Bigl(1 - \frac{1}{n^2}\Bigr),
\qquad
\sum_{n\geq0} \frac{3^n + 4^n}{5^n} .
\]
\end{exercise}

\begin{exercise}[$\star$]\label{exo:b1:series:2}
प्रकृति: $\;\sum \dfrac{n^2}{2^n}$; $\;\sum \dfrac{n!}{n^n}$; $\;\sum
\dfrac{2^n\,n!}{n^n}$; $\;\sum \dfrac{3^n\,n!}{n^n}$। \emph{(\omterm{thm:b1:series:ratio}{अनुपात कसौटी};
$\bigl(1 + \frac1n\bigr)^n \to \eu$ स्मरण कीजिए।)}
\end{exercise}

\begin{exercise}[$\star$]\label{exo:b1:series:3}
प्रकृति: $\;\sum \sin\dfrac{1}{n^2}$; $\;\sum \Bigl(1 - \cos\dfrac1n\Bigr)$;
$\;\sum \dfrac{1}{\sqrt{n(n+1)}}$; $\;\sum \dfrac{\ln n}{n^2}$
\emph{($n^{-3/2}$ से तुलना कीजिए)}।
\end{exercise}

\begin{exercise}[$\star$]\label{exo:b1:series:4}
सिद्ध कीजिए कि $\sum_{n\geq1} \frac{1}{n^2}$ अभिसरित होती है और उसका योग
$\leq 2$ है --- इसके लिए $n \geq 2$ के लिए $\frac{1}{n^2} \leq
\frac{1}{n(n-1)}$ और एक दूरबीनी परिबंध का प्रयोग कीजिए।
\end{exercise}

\begin{exercise}[$\star\star$]\label{exo:b1:series:5}
$\;\sum \dfrac{(-1)^n}{\sqrt n}$ की, $\;\sum \dfrac{(-1)^n}{n + (-1)^n}$ की
\emph{(प्रसार कीजिए: एकांतरित कसौटी सीधे लागू नहीं होती --- क्यों?)} और
$\;\sum \sin\bigl(\pi\sqrt{n^2+1}\,\bigr)$ की \emph{($\pi$ के सापेक्ष घटाइए:
$\sqrt{n^2+1} = n + \frac{1}{2n} + O(n^{-3})$)} प्रकृति बताइए।
\end{exercise}

\begin{exercise}[$\star\star$]\label{exo:b1:series:6}
किन $\alpha > 0$ के लिए $\sum_{n \geq 2} \dfrac{1}{n (\ln n)^{\alpha}}$
अभिसरित होती है? \emph{(समाकल-तुलना; $u = \ln t$ प्रतिस्थापित कीजिए।)}
\end{exercise}

\begin{exercise}[$\star\star$]\label{exo:b1:series:7}
मान लीजिए $u_n = \dfrac{1}{n} - \ln\Bigl(1 + \dfrac1n\Bigr)$। सिद्ध कीजिए कि
$0 \leq u_n \leq \dfrac{1}{2n^2}$, कि $\sum u_n$ अभिसरित होती है, और उससे
ऑयलर के अचर\index{ऑयलर का अचर} का अस्तित्व निकालिए:
\[
\gamma = \lim_{N \to \infty} \Bigl( \sum_{n=1}^{N} \frac 1n - \ln N
\Bigr) .
\]
\end{exercise}

\begin{exercise}[$\star\star$]\label{exo:b1:series:8}
योग संगणित कीजिए
\[
\sum_{n=1}^{\infty} \frac{1}{n(n+2)}
\qquad\text{और}\qquad
\sum_{n=0}^{\infty} \frac{n}{2^n} .
\]
\emph{(पहले के लिए: आंशिक भिन्नें। दूसरे के लिए: $\sum_{n=1}^{N} n x^{n-1}$
संवृत रूप में संगणित कीजिए और $x = \frac12$ पर $N \to \infty$ लीजिए।)}
\end{exercise}

\begin{exercise}[$\star\star\star$]\label{exo:b1:series:9}
(कोशी संघनन) मान लीजिए $(u_n)$ ऋणेतर और ह्रासमान है। सिद्ध कीजिए कि
\[
\sum_{n \geq 1} u_n \text{ अभिसरित होती है}
\iff
\sum_{k \geq 0} 2^k\, u_{2^k} \text{ अभिसरित होती है},
\]
इसके लिए $2$ की क्रमागत घातों के बीच के पदों के पैकेट की तुलना कीजिए। उससे
रीमान कसौटी और \cref{exo:b1:series:6} वापस पाइए।
\end{exercise}

\begin{exercise}[$\star\star\star$]\label{exo:b1:series:10}
$\frac{1}{1+t^2}$ के लिए ढाले हुए \cref{ex:b1:series:ln2} की समाकल-सर्वसमिका
का प्रयोग करके लाइब्निज़ का सूत्र सिद्ध कीजिए
\[
\frac{\pi}{4} = \sum_{n=0}^{\infty} \frac{(-1)^n}{2n+1}
= 1 - \frac13 + \frac15 - \frac17 + \cdots
\]
और त्रुटि-परिबंध $\abs{R_N} \leq \frac{1}{2N+3}$ भी।
\end{exercise}

\begin{exercise}[$\star\star$]\label{exo:b1:series:11}
$\displaystyle\sum_{n \geq 1} \frac{1}{n^{1 + 1/n}}$ की प्रकृति बताइए।
\emph{($n^{1/n}$ की सीमा संगणित कीजिए और व्यापक पद की तुल्यता खोजिए: रीमान
कसौटी को \emph{नियत} घातांक चाहिए।)}
\end{exercise}

\begin{exercise}[$\star\star\star$]\label{exo:b1:series:12}
मान लीजिए $(u_n)$ ऋणेतर और \emph{ह्रासमान} है और $\sum u_n$ अभिसारी है।
सिद्ध कीजिए कि $n\,u_n \to 0$ \emph{($n\,u_{2n}$ को किसी टुकड़े
$\sum_{k=n+1}^{2n} u_k$ से परिबद्ध कीजिए और कोशी कसौटी का प्रयोग कीजिए)}।
दिखाइए कि विलोम विफल है, और यह कि एकदिष्टता की परिकल्पना हटाई नहीं जा सकती।
\end{exercise}

\section{समस्या: ऑयलर का अचर और वह श्रेणी जो अपना योग बदल देती है}

\begin{problem}[{सप्ताहांत समस्या --- $H_n = \ln n + \gamma +
\frac{1}{2n} + O(n^{-2})$, और $1 - \frac12 +
\frac13 - \dots$ को पुनर्क्रमित करके $\frac{\ln 2}{2}$}]
\label{pb:b1:series:1}
हरात्मक \omterm{def:b1:series:def}{श्रेणी} में दो कथाएँ साझा हैं। पहली, उसके अपसरण का यथार्थ लेखा-जोखा:
$H_n - \ln n$ ऑयलर के अचर $\gamma$ की ओर अभिसरित होता है
(\cref{exo:b1:series:7}), और यह समस्या उस \omterm{def:b1:logic:statement}{कथन} को दुतरफ़ा नियम
$\frac{1}{2(n+1)} \leq H_n - \ln n - \gamma \leq \frac{1}{2n}$ तक तीक्ष्ण कर
देती है, तथा $\gamma = 0.5772\dots$ हाथ से प्रमाणित कर देती है। दूसरी,
सप्रतिबंध अभिसरण का कांड: एकांतरित हरात्मक \omterm{def:b1:series:def}{श्रेणी} का योग $\ln 2$ है
(\cref{ex:b1:series:ln2}), फिर भी \emph{वही पद, दूसरे क्रम में}, $\frac{\ln
2}{2}$ देते हैं --- अथवा किसी भी $p, q$ के लिए $\ln 2 + \frac12\ln\frac pq$,
अथवा कोई भी वास्तविक संख्या (रीमान)। दोनों कथाएँ एक ही हैं: पुनर्क्रमित योग
$\gamma$-नियम \emph{के साथ} संगणित किए जाते हैं।\index{ऑयलर का
अचर}\index{श्रेणी का पुनर्क्रमण}

\medskip \noindent\textbf{भाग I --- $\gamma$, घेराव सहित।} $a_n = H_n - \ln
n$ और $b_n = H_n - \ln(n+1)$ रखिए।
\begin{enumerate}
  \item $\frac{t}{1+t} \leq \ln(1 + t) \leq t$ का प्रयोग करके दिखाइए कि
  $(a_n)$ घटता है, $(b_n)$ बढ़ता है, और वे संलग्न हैं; उनकी उभयनिष्ठ सीमा
  $\gamma$ है, और प्रत्येक $n$ के लिए $b_n \leq \gamma \leq a_n$।
  \item पहली संख्यात्मक चाल: $H_{10} = 2.928968\dots$ से $\gamma$ को $b_{10}
  = 0.5311$ और $a_{10} = 0.6264$ के बीच घेरिए। चार दशमलव के लिए इस स्थूल
  घेराव को कितना बड़ा $n$ चाहिए होगा?
  \item यथार्थ पूँछ-निरूपण $a_n - \gamma = \sum_{k \geq n} w_k$ (आंशिक योगों
  की सीमा) दिखाइए, जहाँ
        \[
        w_k = a_k - a_{k+1}
        = \ln\Bigl(1 + \frac1k\Bigr) - \frac{1}{k+1}
        = \int_k^{k+1} \frac{(k + 1 - t)}{t\,(k+1)}\,\dd t ,
        \]
        और समाकल-रूप से दुतरफ़ा परिबंध $\dfrac{1}{2(k+1)^2} \leq w_k \leq
        \dfrac{1}{2k(k+1)}$ निकालिए।
\end{enumerate}

\medskip \noindent\textbf{भाग II --- $\frac{1}{2n}$ नियम।}
\begin{enumerate}[resume]
  \item प्रश्न 3 के परिबंधों का योग कीजिए (दोनों पक्ष दूरबीनी हो जाते हैं या
  दूरबीनों से तुलनीय हैं) और नियम पर पहुँचिए:
        \[
        \frac{1}{2(n+1)} \;\leq\; H_n - \ln n - \gamma \;\leq\;
        \frac{1}{2n} \qquad (n \geq 1).
        \]
  \item $H_n = \ln n + \gamma + \frac{1}{2n} + O\bigl(\frac{1}{n^2}\bigr)$
  निकालिए; ठीक-ठीक कहें तो दिखाइए कि $\gamma_n = H_n - \ln n - \frac{1}{2n}$
  $-\frac{1}{2n(n+1)} \leq \gamma_n - \gamma \leq 0$ संतुष्ट करता है।
  \item $n = 100$ से चार दशमलव प्रमाणित कीजिए: $H_{100} = 5.1873775\dots$
  दिया हो, तो $\gamma_{100} = 0.577207\dots$ संगणित कीजिए और $\gamma =
  0.5772 \pm 5\cdot10^{-5}$ निष्कर्ष निकालिए (सच्चा मान $0.5772156\dots$)।
  \item नियम के दो लाभांश, दोनों आगे आवश्यक: जब $m \to \infty$,
        \[
        H_{2m} - H_m = \ln 2 - \frac{1}{4m} +
        O\Bigl(\frac{1}{m^2}\Bigr),
        \qquad
        \sum_{j=1}^{m} \frac{1}{2j-1} = \frac{\ln m}{2} + \ln 2
        + \frac\gamma2 + o(1) ,
        \]
        जहाँ दूसरा $\sum_{j \leq m} \frac{1}{2j-1} = H_{2m} - \frac12 H_m$
        के द्वारा, और इसी प्रकार $\sum_{j=1}^{m} \frac{1}{2j} = \frac{\ln
        m}{2} + \frac\gamma2 + o(1)$।
\end{enumerate}

\medskip \noindent\textbf{भाग III --- एकांतरित हरात्मक \omterm{def:b1:series:def}{श्रेणी}, द्वितीय कोटि
तक।}
\begin{enumerate}[resume]
  \item (आगमन से, अथवा समूहन से) सर्वसमिका $\sum_{k=1}^{2m}
  \frac{(-1)^{k-1}}{k} = H_{2m} - H_m$ दिखाइए, और उससे योग $\ln 2$ (फिर से)
  तथा यथार्थ गति दोनों निकालिए:
        \[
        \sum_{k=1}^{2m} \frac{(-1)^{k-1}}{k}
        = \ln 2 - \frac{1}{4m} + O\Bigl(\frac{1}{m^2}\Bigr) .
        \]
  \item उससे \emph{किसी भी} सूचकांक पर एकांतरित हरात्मक \omterm{def:b1:series:def}{श्रेणी} की
  अनंतस्पर्शी त्रुटि निकालिए: $S - S_N \sim \frac{(-1)^N}{2N}$ --- अर्थात्
  \cref{thm:b1:series:alternating} के अधिकतम-स्थिति परिबंध $a_{N+1} \approx
  \frac1N$ से दुगुनी छोटी।
  \item (मुफ़्त में त्वरण) दिखाइए कि औसतित योग $\tilde S_N = \frac{S_N +
  S_{N+1}}{2}$ $\tilde S_N = \ln 2 + O\bigl(\frac{1}{N^2}\bigr)$ संतुष्ट
  करते हैं। जाँच: $\ln 2 = 0.69315$ के सामने $S_{10} = 0.64563$, $S_{11} =
  0.73654$, $\tilde S_{10} = 0.69109$: एक औसत लेने से दो दशमलव स्थान मिल
  जाते हैं।
  \item दो वाक्यों में समझाइए कि प्रश्न 2 के घेराव जैसी किसी \emph{धनात्मक}
  अपसारी-पूँछ वाली परिघटना में ऐसी कोई युक्ति क्यों काम नहीं आ सकती:
  एकांतरित त्रुटि दोलन करती है (चिह्न $(-1)^N$), अतः औसत लेने से उसका अग्र
  पद कट जाता है, जबकि $\gamma$-घेराव की त्रुटि $\frac{1}{2n}$ का चिह्न अचर
  है। ($a_n$ और $b_n$ का औसत लेना \emph{सचमुच} काम आता है: $\frac{a_n +
  b_n}{2}$ को मध्यबिंदु आकलन $H_n - \ln\bigl(n + \frac12\bigr)$ से जोड़िए और
  दिखाइए कि उसकी त्रुटि $O\bigl(\frac{1}{n^2}\bigr)$ है।)
\end{enumerate}

\medskip \noindent\textbf{भाग IV --- दृढ़ता और उसकी विफलता।}
\begin{enumerate}[resume]
  \item दिखाइए कि एकांतरित हरात्मक \omterm{def:b1:series:def}{श्रेणी} का धनात्मक भाग $\sum
  \frac{1}{2j-1}$ और ऋणात्मक भाग $\sum \frac{1}{2j}$ दोनों अपसरित होते हैं
  --- यही \emph{सप्रतिबंध} अभिसरण की पहचान है।
  \item प्रश्न 12 के पीछे का व्यापक \omterm{def:b1:logic:statement}{कथन} सिद्ध कीजिए: यदि $\sum u_n$ अभिसरित
  होती है पर $\sum \abs{u_n}$ अपसरित, तो धनात्मक भागों की \omterm{def:b1:series:def}{श्रेणी} $\sum
  u_n^+$ और ऋणात्मक भागों की $\sum u_n^-$ दोनों अपसरित होती हैं
  \emph{($u_n^\pm = \frac{\abs{u_n} \pm u_n}{2}$ से: यदि एक अभिसरित होती, तो
  दूसरी भी होती, अतः $\sum\abs{u_n}$)}। धनात्मक और ऋणात्मक द्रव्यमान का यही
  अक्षय भंडार है जिसे रीमान की विधि ख़र्च करेगी।
  \item (दृढ़ता) सिद्ध कीजिए: यदि $\sum u_n$ \emph{निरपेक्ष रूप से} अभिसरित
  होती है और $\sigma \colon \N \to \N$ एकैकी आच्छादन है, तो $\sum
  u_{\sigma(n)}$ उसी योग की ओर अभिसरित होती है \emph{(बड़े $N$ के लिए पहले
  $M$ पुनर्क्रमित पदों में $u_0, \dots, u_N$ आ जाते हैं; आंशिक योगों की
  तुलना पूँछ $\sum_{n > N}\abs{u_n}$ के द्वारा कीजिए)}।
  \item (रीमान की विधि) मान लीजिए $t \in \R$। एकांतरित हरात्मक \omterm{def:b1:series:def}{श्रेणी} के
  लालची पुनर्क्रमण का वर्णन कीजिए: धनात्मक पद $1, \frac13, \frac15, \dots$
  तब तक लेते जाइए जब तक आंशिक योग पहली बार $t$ पार न कर ले, फिर ऋणात्मक पद
  $-\frac12, -\frac14, \dots$ तब तक जब तक वह पहली बार $t$ से नीचे न गिर जाए,
  और यही दोहराइए। दिखाइए कि हर पद ठीक एक बार प्रयुक्त होता है, कि पहली बार
  पार करने के बाद आंशिक योग $t$ के अंतिम प्रयुक्त पद के भीतर रहते हैं, और
  निष्कर्ष निकालिए कि पुनर्क्रमित \omterm{def:b1:series:def}{श्रेणी} $t$ की ओर अभिसरित होती है: अर्थात्
  \emph{कोई भी} निर्धारित योग प्राप्य है।
\end{enumerate}

\medskip \noindent\textbf{भाग V --- $(p, q)$ सूत्र।} पूर्णांक $p, q \geq 1$
नियत कीजिए। एकांतरित हरात्मक \omterm{def:b1:series:def}{श्रेणी} को खंडों में पुनर्क्रमित कीजिए: $p$
धनात्मक पद (अगले विषम व्युत्क्रम), फिर $q$ ऋणात्मक पद (अगले सम व्युत्क्रम),
और यही दोहराइए।
\begin{enumerate}[resume]
  \item जाँचिए कि यह सच्चा पुनर्क्रमण है (हर पद ठीक एक बार), और यह कि $(p,
  q) = (1, 2)$ के लिए वह यह कहता है
        \[
        1 - \frac12 - \frac14 + \frac13 - \frac16 - \frac18 +
        \frac15 - \frac1{10} - \frac1{12} + \dots
        \]
  \item (यथार्थ आधा करना) $(p, q) = (1, 2)$ के लिए खंड-सर्वसमिका सिद्ध कीजिए
        \[
        \frac{1}{2k-1} - \frac{1}{4k-2} - \frac{1}{4k}
        = \frac12\Bigl(\frac{1}{2k-1} - \frac{1}{2k}\Bigr),
        \]
        और पुनर्क्रमित आंशिक योगों तथा मूल योगों के बीच \emph{यथार्थ} संबंध
        $T_{3K} = \frac12 S_{2K}$ निकालिए: योग का आधा हो जाना हर परिमित चरण
        पर दिखाई देता है, केवल सीमा में नहीं।
  \item दिखाइए कि $K$ पूरे खंडों के बाद आंशिक योग $\sum_{j=1}^{pK}
  \frac{1}{2j-1} - \sum_{j=1}^{qK} \frac{1}{2j}$ के बराबर है, और प्रश्न 7 के
  साथ उसकी सीमा संगणित कीजिए:
        \[
        \ln 2 + \frac12 \ln\frac pq .
        \]
  \item किसी खंड के \emph{भीतर} आंशिक योगों को नियंत्रित कीजिए (पद $0$ की ओर
  जाते हैं) और निष्कर्ष निकालिए कि $(p, q)$-पुनर्क्रमित \omterm{def:b1:series:def}{श्रेणी} $\ln 2 +
  \frac12\ln \frac pq$ की ओर अभिसरित होती है। विशेष रूप से $(1, 2)$
  $\frac{\ln 2}{2}$ देता है: पहले नौ पदों $T_9 = 0.3083$ के सामने सत्यापित
  कीजिए, जो $0.3466$ की ओर रेंग रहे हैं।
  \item संगति की जाँच और परास: $(1,1)$ $\ln 2$ वापस देता है; $(2,1)$
  $\frac32\ln 2$ देता है; $(p, q)$-खंडों से कौन-से योग प्राप्य हैं, और यह
  गणनीय सूची रीमान के पूरे मेनू (प्रश्न 14) के सामने कैसी है?
\end{enumerate}

\medskip \noindent\textbf{भाग VI --- उपसंहार: $\gamma$ काम पर, और संश्लेषण।}
\begin{enumerate}[resume]
  \item अभिसारी \omterm{def:b1:series:def}{श्रेणी} $\sum_{k\geq1} \bigl(\frac1k -
  \ln\frac{k+1}{k}\bigr)$ (\cref{exo:b1:series:7}) का योग पहचानिए: दिखाइए कि
  वह $\gamma$ के बराबर है।
  \item लक्ष्य $t = 1$ के लिए रीमान की विधि (प्रश्न 14) चलाइए और उत्पन्न
  पहले बारह पद ($1, \frac13, -\frac12, \frac15, -\frac14, \frac17, \frac19,
  -\frac16, \frac1{11}, \frac1{13}, -\frac18, \frac1{15}$) सूचीबद्ध कीजिए,
  तथा आंशिक योग ($\approx 0.980$) संगणित कीजिए --- और देखिए कि कलनविधि अपने
  लक्ष्य के इर्द-गिर्द कैसे साँस लेती है।
  \item दिखाइए कि एकांतरित हरात्मक \omterm{def:b1:series:def}{श्रेणी} का कोई पुनर्क्रमण $+\infty$ की ओर
  अपसरित होता है \emph{(धनात्मक पदों के इतने लंबे खंड कि हर बार $1$ का लाभ
  हो, प्रश्न 12 का प्रयोग करते हुए, और उनके बीच एक-एक ऋणात्मक पद)}।
  \item $\gamma$-नियम से \cref{ex:b1:series:harmonicstack} को तीक्ष्ण कीजिए:
  दिखाइए कि $H_N \geq 20$ वाला पहला सूचकांक $N = \eu^{\,20 - \gamma}\,(1 +
  o(1)) \approx 2.7\cdot10^{8}$ संतुष्ट करता है --- अर्थात् \omterm{pb:b1:series:1}{ऑयलर का अचर} ठीक
  वही सुधार है जो स्थूल घेराव से छूट रहा था।
  \item संश्लेषण, एक-एक वाक्य में: (क) $\gamma$-नियम और उसके तीनों अंश ($\ln
  n$, $\gamma$, $\frac{1}{2n}$) क्या-क्या देते हैं; (ख) सप्रतिबंध अभिसरण योग
  को क्रम-निर्भर क्यों बना देता है जबकि \omterm{thm:b1:series:absolute}{निरपेक्ष अभिसरण} उसे मना कर देता है;
  (ग) $(p,q)$ सूत्र किसी अमूर्त अस्तित्व-दावे के बदले $\gamma$-नियम के साथ
  एक \emph{संगणना} कैसे था; (घ) ये धागे कहाँ आगे चलते हैं --- स्नातक वर्ष 2
  के खंड में घात \omterm{def:b1:series:def}{श्रेणी} और श्रेणियों के गुणनफल, तथा स्नातक वर्ष 3 के खंड की
  स्टर्लिंग सूत्र वाली सप्ताहांत समस्या, जहाँ वही योग-बनाम-समाकल लेखा-जोखा
  पूरी शक्ति से चलता है।
\end{enumerate}
\end{problem}
